% Part: first-order-logic
% Chapter: beyond
% Section: many-sorted-logic

\documentclass[../../../include/open-logic-section]{subfiles}

\begin{document}

\olfileid{fol}{byd}{msl}

\olsection{Logika Akeh Jinis}

Ing logika sepisanan, variabel lan kuantor
nglimputi mung siji !!{domain}. Nanging asring
migunani yen ana maneka !!{domain}s
(kang padha pisah): upamane, kita bisa mbutuhake
!!a{domain} wilangan, !!a{domain} objek geometri,
!!a{domain} fungsi saka wilangan menyang
wilangan, !!a{domain} grup abelian, lan liya-liyane.

Logika akeh jinis nyedhiyakake kerangka kaya
mangkono. Kita miwiti saka dhaptar "jinis";
"jinis" sawijining objek nuduhake
"!!{domain}" kang kudune dienggoni objek iku.
Sabanjure ana !!{variable}s lan kuantor
kanggo saben jinis, lan (lumrahe) siji
!!{identity} kanggo saben jinis. Fungsi
lan relasi uga diwatesi miturut "jinis"
objek kang bisa dadi argumene. Saliyane iku,
aturan logika sepisanan lumrahe tetep
digunakake, kanthi aturan kuantor kang
diulang kanggo saben jinis.

Upamane, kanggo nliti hubungan internasional
kita bisa milih basa kanthi rong jinis objek,
warga Prancis lan warga Jerman. Kita bisa
nduweni relasi unari "ngombe anggur" kanggo
objek jinis kapisan; relasi unari liyane
"mangan wurst" kanggo objek jinis kapindho;
lan relasi binari "mbentuk pasangan bojo
lintas negara" kanthi rong argumen, kang
argumen kapisane saka jinis kapisan lan
argumen kapindhone saka jinis kapindho.
Yen kita nggunakake variabel $a$, $b$, $c$
kanggo warga Prancis lan $x$, $y$, $z$
kanggo warga Jerman, mula
\[
\lforall[a][\lforall x][(\Atom{\Obj{MarriedTo}}{a,x} \lif
(\Atom{\Obj{DrinksWine}}{a} \lor \lnot \Atom{\Obj{EatsWurst}}{x})]]
\]
nyatakake yen sembarang warga Prancis nikah
karo warga Jerman, mula warga Prancis iku
ngombe anggur utawa warga Jerman iku ora
mangan wurst.

Logika akeh jinis bisa dilebokake ing logika
sepisanan kanthi cara kang lumrah: kabeh
objek saka !!{domain}s kang akeh jinis iku
digabung dadi siji !!{domain} logika sepisanan,
!!{predicate}s unari digunakake kanggo
nyathet jinis-jinis kasebut, lan kuantor
diwatesi miturut jinis. Upamane, basa
logika sepisanan kang cocog karo tuladha
ing ndhuwur bakal nduweni !!{predicate}s
unari "$\Obj{German}$" lan "$\Obj{French}$",
ing sajabane relasi liyane kang wis kasebut,
kanthi syarat jinis ing relasi kasebut dicopot.
Kuantor kang diwatesi jinis
$\lforall[x][!A]$, ing ngendi $x$ iku
!!a{variable} saka jinis Jerman, diterjemahake dadi
\[
\lforall[x][(\Atom{\Obj{German}}{x} \lif !A)].
\]
Kita kudu nambah aksioma kang njamin
yen jinis-jinis iku padha pisah---upamane,
$\lforall[x][\lnot (\Atom{\Obj{German}}{x} \land
  \Atom{\Obj{French}}{x})]$---uga aksioma
kang njamin yen "ngombe anggur" mung bener
kanggo objek kang nyukupi predikat
$\Atom{\Obj{French}}{x}$, lan liya-liyane.
Kanthi tata cara lan aksioma kasebut,
ora angel nuduhake yen !!{sentence}s
logika akeh jinis bisa diterjemahake
dadi !!{sentence}s logika sepisanan,
lan !!{derivation}s logika akeh jinis
bisa diterjemahake dadi !!{derivation}s
logika sepisanan. Kajaba iku,
!!{structure}s logika akeh jinis uga
"diterjemahake" dadi !!{structure}s
logika sepisanan kang cocog lan
kosok balene, mula kita uga nduweni
teorema kelengkapan kanggo logika
akeh jinis.

\end{document}
